\documentclass[12pt,notitlepage]{report}
% The shared notes preamble lives outside this repo.
\input{/Users/aryankulkarni/Documents/y3-notes/preamble.tex}

\title{Week 2: exploring the data and a first linear model}
\author{Aryan Kulkarni}
\date{MTH6134 group project}

\begin{document}
\maketitle

\paragraph{Findings.}
Playing minutes depend on position and training load and on nothing else we
measure: the model explains 28\% of the variation and passes every check of
the Week 2 practical. Sixteen players sit on a ceiling of 3960 minutes, which
week 10 takes up. BMI shows no association with any other variable.

\section{The data}

\begin{minted}{r}
source("../common.R")

out(glimpse(players), "glimpse-players")
out(summary(players), "summary-players")
out(glimpse(matches), "glimpse-matches")
\end{minted}

\inputminted{text}{output/glimpse-players.txt}
\inputminted{text}{output/summary-players.txt}
\inputminted{text}{output/glimpse-matches.txt}

\mintinline{r}|common.R| reads both files, fixes the factor levels (DEF and
Away are the reference levels) and derives \mintinline{r}|opponent|,
\mintinline{r}|goal_diff| and \mintinline{r}|result| for the matches.

\begin{minted}{r}
checks <- c(
  missing_values = sum(is.na(players)) + sum(is.na(matches)),
  clubs = nlevels(players$team),
  smallest_squad = min(table(players$team)),
  largest_squad = max(table(players$team)),
  missing_rank = setdiff(1:24, matches$opponent_rank),
  players_at_3960 = sum(players$minutes_played == 3960)
)
out(checks, "checks")
out(table(matches$home_away, matches$result), "results")
\end{minted}

\inputminted{text}{output/checks.txt}
\inputminted{text}{output/results.txt}

Nothing is missing. The 24 squads hold 24 to 28 players. Ashford meet each of
23 opponents home and away, and the one rank absent from
\mintinline{r}|opponent_rank| is 9, so Ashford finished ninth (22 wins, 11
draws, 13 defeats, 77 points). The largest value of
\mintinline{r}|minutes_played| is 3960, which is 44 full matches, and 16
players share it, 12 of them goalkeepers. That pattern points to a recording
ceiling.

\paragraph{Variable types.}
\mintinline{r}|read_csv| only decides how a column is stored, so the types
below are our reading of what each variable means.

\begin{center}
  \begin{tabular}{lll}
    \toprule
    Data set & Type & Variables \\
    \midrule
    players & nominal          & \mintinline{r}|team|, \mintinline{r}|position| \\
    players & binary           & \mintinline{r}|injured| \\
    players & ordinal (1 to 5) & \mintinline{r}|training_load| \\
    players & count            & \mintinline{r}|prev_injuries| \\
    players & continuous       & \mintinline{r}|age|, \mintinline{r}|bmi|, \mintinline{r}|minutes_played| \\
    matches & binary           & \mintinline{r}|home_away| \\
    matches & count            & \mintinline{r}|goals_scored|, \mintinline{r}|goals_conceded| \\
    matches & rank             & \mintinline{r}|opponent_rank| \\
    matches & ordinal, derived & \mintinline{r}|result| \\
    \bottomrule
  \end{tabular}
\end{center}

\section{Plots before modelling}

\begin{minted}{r}
fig <- players |>
  pivot_longer(c(minutes_played, bmi, age)) |>
  ggplot(aes(x = value)) +
  geom_histogram(bins = 20, colour = "white") +
  facet_wrap(~name, scales = "free") +
  labs(x = NULL, y = "Players")
save_fig(fig, "distributions", height = 2.3)

fig <- players |>
  mutate(load = factor(training_load), previous = factor(prev_injuries)) |>
  pivot_longer(
    c(position, load, previous),
    names_to = "variable", values_to = "level",
    values_transform = as.character
  ) |>
  mutate(
    variable = factor(variable, levels = c("position", "load", "previous")),
    level = factor(level, levels = c(levels(players$position), 0:5))
  ) |>
  ggplot(aes(x = level, y = minutes_played)) +
  geom_boxplot(outlier.size = 0.6) +
  facet_grid(
    ~variable, scales = "free_x", space = "free_x",
    labeller = as_labeller(c(
      position = "position", load = "training load",
      previous = "previous injuries"
    ))
  ) +
  labs(x = NULL, y = "Minutes played")
save_fig(fig, "minutes-by-group", height = 2.6)

fig <- players |>
  pivot_longer(c(age, bmi)) |>
  ggplot(aes(x = value, y = minutes_played)) +
  geom_point(alpha = 0.3) +
  geom_smooth(method = "loess", formula = y ~ x) +
  facet_wrap(~name, scales = "free_x") +
  labs(x = NULL, y = "Minutes played")
save_fig(fig, "minutes-by-number", height = 2.4)
\end{minted}

\begin{figure}[H]
  \centering
  \incfig{distributions}
  \caption{The three continuous variables are roughly symmetric with no long
    tail, so a normal error model is plausible for either outcome.}
\end{figure}

\begin{figure}[H]
  \centering
  \incfig{minutes-by-group}
  \caption{Minutes differ by position (goalkeepers most, forwards least) and
    rise evenly across training load. Previous injuries show no pattern.}
  \label{fig:groups}
\end{figure}

\begin{figure}[H]
  \centering
  \incfig{minutes-by-number}
  \caption{The smooth is flat for age and for BMI: neither carries a trend,
    and neither bends.}
\end{figure}

\section{A linear model for minutes played}

For player $i$ the model is
\[
  Y_i = \beta_0 + \bm{x}_i^\top \bm{\beta} + \varepsilon_i,
  \qquad \varepsilon_i \iid \Norm(0, \sigma^2),
\]
which assumes a mean that is linear in the explanatory variables,
independent errors, one error variance and normal errors.
\mintinline{r}|injured| is left out because an injury during the season
removes playing minutes, so it is a consequence of the same season and cannot
explain it. \mintinline{r}|team| waits for week 11.

\begin{minted}{r}
fit_full <- lm(
  minutes_played ~ position + age + training_load + prev_injuries + bmi,
  data = players
)
out(coef(summary(fit_full)), "coef-fit-full")

fit_min <- lm(minutes_played ~ position + training_load, data = players)
fit_factor <- lm(minutes_played ~ position + factor(training_load), players)
fit_inter <- lm(minutes_played ~ position * training_load, players)
tests <- rbind(
  "drop age, prev_injuries, bmi" = anova(fit_min, fit_full)[2, 3:6],
  "load as a factor" = anova(fit_min, fit_factor)[2, 3:6],
  "position:load interaction" = anova(fit_min, fit_inter)[2, 3:6]
)
out(tests, "tests")
\end{minted}

\inputminted{text}{output/coef-fit-full.txt}
\inputminted{text}{output/tests.txt}

Age, previous injuries and BMI have $t$ values below 1 in absolute size, and
the $F$ test for dropping all three gives $p = 0.91$. Following the practical,
two more $F$ tests settle the form of what remains. Treating training load as
a five-level factor does not improve on a straight line ($p = 0.75$), and the
slope does not differ by position ($p = 0.18$). The model is therefore
additive in position and training load.

\begin{minted}{r}
out(summary(fit_min), "summary-fit-min")
out(confint(fit_min), "confint-fit-min")
\end{minted}

\inputminted{text}{output/summary-fit-min.txt}
\inputminted{text}{output/confint-fit-min.txt}

The fitted model for a defender is
$\widehat{\E}(\text{minutes}) = 2266 + 144 \times \text{load}$. Each extra
point of training load goes with 144 more minutes (95\% interval 96 to 191),
about a match and a half. At equal load a midfielder plays 322 minutes fewer
than a defender, a forward 622 fewer and a goalkeeper 578 more. This is an
association: coaches may train their regular starters harder, so the slope
does not say that raising a player's load earns him minutes.

$R^2 = 0.28$ and the residual standard error is 604 minutes, close to seven
full matches, so the model describes group averages well and individual
players poorly.

\section{Checking the assumptions}

The checks are those of the practical: studentised residuals against fitted
values, their normal QQ plot, Cook's distances against the threshold
$8/(n - 2p)$, and the Box-Cox profile likelihood.

\begin{minted}{r}
r <- rstudent(fit_min)
cook <- cooks.distance(fit_min)
n <- nobs(fit_min)
p <- length(coef(fit_min))
cook_cut <- 8 / (n - 2 * p)
flagged <- which(cook > cook_cut)

base_fig("checks-min", mfrow = c(2, 2), height = 4.6)
plot(fitted(fit_min), r, col = pt,
     xlab = "Fitted values", ylab = "Studentised residuals")
abline(h = 0, lwd = 2)
qqnorm(r, main = "", col = pt)
qqline(r, lwd = 2)
plot(cook, type = "h", col = blue, ylim = c(0, 1.15 * max(cook)),
     xlab = "Observation", ylab = "Cook's distance")
abline(h = cook_cut, lty = 2)
text(flagged, cook[flagged], labels = flagged, pos = 3, cex = 0.8)
bc <- MASS::boxcox(fit_min, lambda = seq(0, 2, by = 0.01))
dev.off()

lambda_ci <- range(bc$x[bc$y > max(bc$y) - qchisq(0.95, 1) / 2])
out(c(lambda = bc$x[which.max(bc$y)], lower = lambda_ci[1],
      upper = lambda_ci[2], largest_abs_residual = max(abs(r)),
      cook_threshold = cook_cut), "check-numbers")

influential <- players[flagged, c("position", "minutes_played", "training_load")] |>
  as.data.frame() |>
  mutate(fitted = round(fitted(fit_min)[flagged]),
         rstudent = round(r[flagged], 2), cook = round(cook[flagged], 3))
rownames(influential) <- flagged
out(influential, "influential")

fit_drop <- update(fit_min, data = players[-flagged, ])
out(round(cbind(all = coef(fit_min), without = coef(fit_drop)), 1),
    "sensitivity")
\end{minted}

\begin{figure}[H]
  \centering
  \incfig{checks-min}
  \caption{Checks on the final model. The dashed line in the Cook's distance
    panel is $8/(n - 2p)$.}
  \label{fig:checks}
\end{figure}

\inputminted{text}{output/check-numbers.txt}
\inputminted{text}{output/influential.txt}
\inputminted{text}{output/sensitivity.txt}

\begin{itemize}
  \item \emph{Residuals against fitted values.} The points scatter evenly
    about zero with no curve and no funnel, so linearity and constant
    variance raise no concern. No studentised residual exceeds 3 in absolute
    size.
  \item \emph{QQ plot.} The points follow the line, so normality is
    reasonable. The upper tail falls slightly below the line, which is the
    ceiling at 3960 holding down the largest residuals.
  \item \emph{Cook's distance.} Four players exceed the threshold of 0.013,
    each with far fewer or far more minutes than his position and load
    predict. The largest distance is 0.028, which is small in absolute terms.
    Refitting without the four leaves every conclusion unchanged: the load
    slope moves from 144 to 130 and the goalkeeper effect from 578 to 611.
    They stay in.
  \item \emph{Box-Cox.} The profile likelihood peaks at $\lambda = 1.15$ and
    its 95\% interval, 0.96 to 1.36, contains 1. No transformation of the
    response is needed.
\end{itemize}

\section{Confidence and prediction intervals}

\begin{minted}{r}
grid <- expand_grid(position = levels(players$position), training_load = 1:5)
conf <- predict(fit_min, grid, interval = "confidence")
pred <- predict(fit_min, grid, interval = "prediction")
bands <- grid |>
  mutate(
    position = factor(position, levels = levels(players$position)),
    fit = conf[, "fit"], conf_lwr = conf[, "lwr"], conf_upr = conf[, "upr"],
    pred_lwr = pred[, "lwr"], pred_upr = pred[, "upr"]
  )

fig <- ggplot(bands, aes(x = training_load)) +
  geom_jitter(aes(y = minutes_played), data = players, colour = grey,
              width = 0.15, height = 0, alpha = 0.25, size = 0.6) +
  geom_hline(yintercept = 3960, linetype = "dotted", colour = grey) +
  geom_ribbon(aes(ymin = conf_lwr, ymax = conf_upr), fill = blue, alpha = 0.35) +
  geom_line(aes(y = fit)) +
  geom_line(aes(y = pred_lwr), linetype = "dashed") +
  geom_line(aes(y = pred_upr), linetype = "dashed") +
  facet_wrap(~position, nrow = 1) +
  labs(x = "Training load", y = "Minutes played")
save_fig(fig, "intervals", height = 2.8)

new <- data.frame(position = c("MID", "GK"), training_load = c(3, 5))
out(cbind(new, round(predict(fit_min, new, interval = "prediction"))),
    "predictions")
\end{minted}

\begin{figure}[H]
  \centering
  \incfig{intervals}
  \caption{Fitted lines with 95\% confidence bands (shaded) and prediction
    intervals (dashed). The dotted line is the ceiling of 3960 minutes.}
  \label{fig:intervals}
\end{figure}

\inputminted{text}{output/predictions.txt}

The confidence bands are narrow because 620 players pin down each mean: a
midfielder at load 3 averages 2375 minutes, give or take 83. The prediction
interval for one such player runs from 1185 to 3565, a width set almost
entirely by $\sigma$. For a goalkeeper at load 5 the upper prediction limit is
4761 minutes, more than the 4140 a 46-match season contains. A normal model
has no upper bound, and this is where the ceiling starts to matter.

\section{BMI as a second outcome}

\begin{minted}{r}
fit_bmi <- lm(
  bmi ~ position + age + training_load + prev_injuries + minutes_played,
  data = players
)
f <- summary(fit_bmi)$fstatistic
out(c(r_squared = summary(fit_bmi)$r.squared, f = f[[1]],
      p_value = pf(f[[1]], f[[2]], f[[3]], lower.tail = FALSE),
      smallest_coefficient_p = min(coef(summary(fit_bmi))[-1, 4])),
    "bmi")
\end{minted}

\inputminted{text}{output/bmi.txt}

The overall $F$ test gives $p = 0.73$, $R^2$ is below 0.01, and no single
coefficient is close to significant. BMI shows no association with position,
age, training load, injury history or playing time. There is no model to build
here, and BMI can be set aside as a response.

\section{Summary}

\begin{itemize}
  \item Minutes played is modelled by position and a linear training-load
    term. The assumptions hold and no transformation is needed.
  \item The model predicts averages, with $R^2 = 0.28$ and a residual standard
    error of 604 minutes.
  \item Carried forward: the ceiling at 3960 minutes (week 10) and the
    grouping of players in clubs, which the independence assumption ignores
    (week 11).
\end{itemize}

\end{document}
